\section{Exact branch cylinders at a retained presentation modulus}

Piecewise-affine iteration carries three different moduli: the retained
presentation modulus, a sufficient word-stability modulus, and the least
period of the actual word cylinder.  They need not coincide.  Keeping all
three visible gives both the coarse coordinate used in computation and the
smallest coordinate justified by a selected word.

\begin{definition}[periodic affine branch system]
\label{def:residue-affine-system}
Fix $M\ge2$.  Let
\[
 \mathbb Z=\coprod_{b\in\mathcal B}X_b
\]
be a finite partition in which every $X_b$ is a union of classes modulo
$M$.  Write $m_b$ for the least positive translation period of $X_b$; thus
$m_b\mid M$.  Choose $a_b,c_b\in\mathbb Z$ and $d_b\in\mathbb N$ with
$d_b\mid m_b$, and suppose that
\[
 F(x)=\frac{a_bx+c_b}{d_b}\in\mathbb Z\qquad(x\in X_b)
 \tag{1.1}\label{eq:branch-system}
\]
defines a map $F:\mathbb Z\to\mathbb Z$.  Zero slopes are allowed.
\end{definition}

For a word $w=(b_0,\ldots,b_{\ell-1})$, retain the unreduced integers
\[
 A_0=D_0=1,\qquad B_0=0,
\]
and define
\begin{equation}
 A_{j+1}=a_{b_j}A_j,\qquad
 D_{j+1}=d_{b_j}D_j,\qquad
 B_{j+1}=a_{b_j}B_j+c_{b_j}D_j.
 \label{eq:affine-word-recurrence}
\end{equation}
The length-$\ell$ word cylinder is
\[
 \mathcal C_w=
 \{x\in\mathbb Z:F^j(x)\in X_{b_j}\text{ for }0\le j<\ell\}.
\]
For $0\le j<\ell$, put
\begin{equation}
 \lambda_{j,w}=
 \frac{m_{b_j}D_j}{\gcd(|A_j|,m_{b_j}D_j)},
 \qquad
 \Lambda_w=\operatorname{lcm}_{0\le j<\ell}\lambda_{j,w}.
 \label{eq:coefficient-stability-modulus}
\end{equation}

\begin{theorem}[sufficient and least branch-cylinder coordinates]
\label{thm:general-branch-cylinder}
For every $x\in\mathcal C_w$ and $0\le j\le\ell$,
\begin{equation}
 F^j(x)=\frac{A_jx+B_j}{D_j}.
 \label{eq:general-affine-word}
\end{equation}
The set $\mathcal C_w$ is invariant under translation by $\Lambda_w$, and
\begin{equation}
 \Lambda_w\mid\prod_{j=0}^{\ell-1}m_{b_j}\mid M^\ell.
 \label{eq:stability-divisibility}
\end{equation}

Let $Q=M^\ell$ and let $R_w\subseteq\mathbb Z/Q\mathbb Z$ be the exact set
of residues whose lifts lie in $\mathcal C_w$.  Its translation stabilizer
has the unique form
\[
 \{u\in\mathbb Z/Q\mathbb Z:R_w+u=R_w\}
   =\pi_w\mathbb Z/Q\mathbb Z
\]
for a unique divisor $\pi_w\mid Q$.  This $\pi_w$ is the least positive
translation period of the cylinder, and
\begin{equation}
 \pi_w\mid\Lambda_w.
 \label{eq:least-period-divides-stability}
\end{equation}
If $\mathcal C_w\ne\varnothing$, then also
\begin{equation}
 \frac{D_\ell}{\gcd(|A_\ell|,D_\ell)}\mid\pi_w.
 \label{eq:integrality-lower-period}
\end{equation}

For every occupied component $r+\pi_w\mathbb Z\subseteq\mathcal C_w$,
\begin{equation}
 F^\ell(r+\pi_wq)=F^\ell(r)+\kappa_wq,
 \qquad
 \kappa_w=\frac{A_\ell\pi_w}{D_\ell}\in\mathbb Z.
 \label{eq:least-cylinder-coordinate}
\end{equation}
If $A_\ell\ne0$, this restriction is a bijection onto
$F^\ell(r)+\kappa_w\mathbb Z$, with inverse
\begin{equation}
 z\longmapsto
 r+\pi_w\frac{z-F^\ell(r)}{\kappa_w},
 \label{eq:component-inverse}
\end{equation}
defined exactly when $\kappa_w\mid z-F^\ell(r)$.  If $A_\ell=0$, the
restriction is constant.  Full fibres of $F^\ell$ are unions of these
component fibres; different components are not presumed disjoint in the
image.
\end{theorem}

\begin{proof}
Induction in \eqref{eq:branch-system} gives
\eqref{eq:general-affine-word} with the unreduced recurrence
\eqref{eq:affine-word-recurrence}.  Suppose $x\in\mathcal C_w$ and
$y-x\in\Lambda_w\mathbb Z$.  If the two inputs have followed $w$ through
time $j$, then
\[
 F^j(y)-F^j(x)=\frac{A_j(y-x)}{D_j}.
\]
The definition of $\lambda_{j,w}$ is exactly the divisibility condition
$m_{b_j}D_j\mid A_j(y-x)$.  Hence the last difference is a multiple of
$m_{b_j}$, so $F^j(y)$ lies in the same branch $X_{b_j}$.  Induction proves
that $y\in\mathcal C_w$ and also proves the difference formula at every
stage.

Since $d_{b_u}\mid m_{b_u}$,
\[
 \lambda_{j,w}\mid m_{b_j}D_j
 \mid\prod_{u=0}^{j}m_{b_u}
 \mid\prod_{u=0}^{\ell-1}m_{b_u}.
\]
Taking the least common multiple proves the first divisibility in
\eqref{eq:stability-divisibility}; the second follows from $m_b\mid M$.
Thus both $Q$ and $\Lambda_w$ are periods of $\mathcal C_w$.

Every subgroup of the cyclic group $\mathbb Z/Q\mathbb Z$ is uniquely
$\pi\mathbb Z/Q\mathbb Z$ for a divisor $\pi$ of $Q$.  Applying this to the
translation stabilizer defines $\pi_w$.  Since the class of $\Lambda_w$
belongs to that stabilizer and $\Lambda_w\mid Q$, one has
$\pi_w\mid\Lambda_w$.  If the cylinder is nonempty, choose
$x\in\mathcal C_w$.  Both $x$ and $x+\pi_w$ follow $w$, and their integral
endpoints differ by $A_\ell\pi_w/D_\ell$.  Consequently
$D_\ell\mid A_\ell\pi_w$, which is equivalent to
\eqref{eq:integrality-lower-period}.  The same subtraction proves
\eqref{eq:least-cylinder-coordinate}.  The inverse and fibre statements are
then the literal solution of the displayed affine equation.
\end{proof}

\begin{corollary}[the retained $M^\ell$ coordinate]
\label{cor:retained-cylinder-coordinate}
If $x\equiv y\pmod{M^\ell}$, then their first $\ell$ branch symbols agree and
\begin{equation}
 F^\ell(x)-F^\ell(y)=\frac{A_\ell}{D_\ell}(x-y).
 \label{eq:general-cylinder-difference}
\end{equation}
In particular, for $n\ge0$, $r=n\bmod M^\ell$, and
$q=(n-r)/M^\ell$,
\begin{equation}
 F^\ell(n)=F^\ell(r)+A_\ell\frac{M^\ell}{D_\ell}q.
 \label{eq:general-cylinder-coordinate}
\end{equation}
\end{corollary}

\begin{proof}
Theorem~\ref{thm:general-branch-cylinder} gives
$\Lambda_w\mid M^\ell$ for the word of either input, so congruence modulo
$M^\ell$ preserves that word.  Subtraction of
\eqref{eq:general-affine-word} gives both formulas.
\end{proof}

\subsection{Exact relation to the generalized-map literature}

When all branch slopes are nonzero, the retained presentation above embeds
directly into Matthews's common-denominator formalism.  Refine each $X_b$
into its individual classes $i\pmod M$ and, on a class contained in $X_b$,
put
\begin{equation}
 \mu_i=\frac{M}{d_b}a_b,
 \qquad
 \rho_i=-\frac{M}{d_b}c_b.
 \label{eq:matthews-embedding}
\end{equation}
Then $\rho_i\equiv i\mu_i\pmod M$ and
\[
 \frac{\mu_i x-\rho_i}{M}=\frac{a_bx+c_b}{d_b}.
\]
This is an exact retained-modulus change of presentation, not a minimization
of the modulus.  Matthews's general affine-iterate formula therefore contains
the same raw word algebra
\cite[author-form pp.~2--4, Eq.~(3), Theorems~2.1--2.2]{Matthews2010}.
Kohl's residue-class-wise affine maps give the corresponding reduced-triple
language \cite[Definition~1.1 and Remark~1.2]{Kohl2008}.

The stronger literature conclusion that each length-$\ell$ word is one class
modulo $M^\ell$ requires the unit hypotheses on every retained multiplier.
The same gate controls Haar preservation and mixing
\cite[pp.~172--174, Lemmas~3--5 and Corollary~1]{MatthewsWatts1984}.
Theorem~\ref{thm:general-branch-cylinder} makes no such inference when a
multiplier is a nonunit.

\subsection{The generalized \texorpdfstring{$5x+1$}{5x+1} map}

Consider Oliveira e Silva's one-branch-per-iterate map
\begin{equation}
 C(n)=
 \begin{cases}
 n/2,&n\equiv0\pmod2,\\
 n/3,&n\not\equiv0\pmod2,\ n\equiv0\pmod3,\\
 5n+1,&\gcd(n,6)=1.
 \end{cases}
 \label{eq:generalized-five-map}
\end{equation}
Its three branch sets have least periods $2,6,6$ and branch triples
$(a,c,d)=(1,0,2),(1,0,3),(5,1,1)$.  At retained modulus $6$, the exact
Matthews pairs $(\mu_i,\rho_i)$ are
\begin{equation}
 (3,0)\ (i=0,2,4),\qquad
 (2,0)\ (i=3),\qquad
 (30,-6)\ (i=1,5).
 \label{eq:five-map-matthews-data}
\end{equation}
All retained multipliers are nonunits modulo $6$.

\begin{corollary}[optimal and universal coordinates for a fixed word]
\label{cor:corrected-six-cylinder}
Let $w$ be a nonempty length-$\ell$ word for $C$, containing $k_2,k_3,k_5$
branches of the three respective types, and suppose $\mathcal C_w$ is
nonempty.  Put
\[
 D_w=2^{k_2}3^{k_3},\qquad A_w=5^{k_5}.
\]
Then
\begin{equation}
 D_w\mid\pi_w\mid\Lambda_w
 \mid 2^{k_2}6^{k_3+k_5}\mid6^\ell.
 \label{eq:five-map-period-chain}
\end{equation}
On every occupied class $r+\pi_w\mathbb Z$,
\begin{equation}
 C^\ell(r+\pi_wq)=C^\ell(r)
       +5^{k_5}\frac{\pi_w}{D_w}q.
 \label{eq:optimal-five-map-coordinate}
\end{equation}
The universal retained coordinate is obtained by writing
$n=r_6+6^\ell q$, where $0\le r_6<6^\ell$:
\begin{equation}
 C^\ell(n)=C^\ell(r_6)
       +5^{k_5}\frac{6^\ell}{D_w}q.
 \label{eq:corrected-six-cylinder}
\end{equation}
The representative $r_6$ is canonical modulo $6^\ell$; it need not be the
canonical remainder modulo $D_w$.
\end{corollary}

\begin{proof}
The product of the least branch periods is
$2^{k_2}6^{k_3+k_5}$, so the upper divisibilities follow from
Theorem~\ref{thm:general-branch-cylinder}.  Since
$\gcd(5^{k_5},D_w)=1$, the lower divisibility
\eqref{eq:integrality-lower-period} is $D_w\mid\pi_w$.  The two coordinate
formulas are \eqref{eq:least-cylinder-coordinate} and
\eqref{eq:general-cylinder-coordinate}.
\end{proof}

\begin{corollary}[inverse enumeration, densities, and coalescence]
\label{cor:five-map-enumeration}
Let $r_1,\ldots,r_s$ be the occupied residues modulo $\pi_w$.  Put
\[
 z_u=C^\ell(r_u),\qquad
 \kappa_w=5^{k_5}\pi_w/D_w>0.
\]
For $z\in\mathbb Z$, all inputs in $\mathcal C_w$ with endpoint $z$ are
exactly
\[
 r_u+\pi_w\frac{z-z_u}{\kappa_w}
 \quad(1\le u\le s,\ \kappa_w\mid z-z_u).
 \tag{1.14}\label{eq:five-map-full-fibre}
\]
Each positive input ray has natural density $1/\pi_w$ and each positive
image ray has natural density $1/\kappa_w$.  These image densities may be
added only after their intersections have been checked.

More generally, image progressions $z+\kappa\mathbb Z$ and
$z'+\kappa'\mathbb Z$ meet if and only if
\begin{equation}
 \gcd(\kappa,\kappa')\mid z'-z.
 \label{eq:two-cylinder-crt}
\end{equation}
When they meet, the Chinese remainder theorem gives one class modulo
$\operatorname{lcm}(\kappa,\kappa')$.  Finally, a positive integer $h$ is a
translation period of the word cylinder if and only if $\pi_w\mid h$.
\end{corollary}

\begin{proof}
Formula \eqref{eq:five-map-full-fibre} is the union of the exact component
inverses \eqref{eq:component-inverse}.  The density assertions are the usual
counts for one arithmetic progression; no disjointness has been assumed.
Condition \eqref{eq:two-cylinder-crt} is the exact solvability criterion for
$z+\kappa q=z'+\kappa'q'$.  The final assertion is the definition of the
least positive generator of the translation stabilizer.
\end{proof}

\begin{proposition}[obstruction to denominator-only reduction]
\label{prop:denominator-obstruction}
For the map \eqref{eq:generalized-five-map}, no one-step coordinate can choose
the branch from the canonical remainder modulo the post-selected denominator
alone.
\end{proposition}

\begin{proof}
Take $n=1$.  Its first branch is $5n+1$, whose denominator product is
$D_1=1$.  The canonical remainder of $1$ modulo $1$ is $0$, but $0$ takes
the division-by-$2$ branch.  Thus reduction modulo the denominator erases the
branch symbol in the smallest possible case.  Equivalently, branch selection
cannot factor through the one-point quotient $\mathbb Z/1\mathbb Z$.
\end{proof}

\begin{nonclaim}
The coordinate theorems prove affine transport, exact component fibres, and
the displayed density facts.  They do not imply boundedness, periodicity, or
convergence of integer orbits.  In particular, the nonunit data
\eqref{eq:five-map-matthews-data} do not inherit the singleton-cylinder or
Haar-mixing conclusions of the Matthews unit theory.
\end{nonclaim}
